% =====================================================================
%  TCET CODING PLATFORM -- Official Technical Documentation
%  MODULE 7 of 7 : Deployment, Infrastructure, Operations & Future Scope
%  STANDALONE FILE.  Compile: pdflatex Module7_Deployment_Future.tex  (twice)
% =====================================================================
\documentclass[12pt,a4paper,oneside]{report}
\usepackage[T1]{fontenc}
\usepackage[utf8]{inputenc}
\usepackage[a4paper,margin=1in]{geometry}
\usepackage{mathptmx}\usepackage{microtype}\usepackage{graphicx}\usepackage{xcolor}
\usepackage{tabularx}\usepackage{booktabs}\usepackage{longtable}\usepackage{xltabular}
\usepackage{array}\usepackage{enumitem}\usepackage{listings}\usepackage{titlesec}
\usepackage{fancyhdr}\usepackage{parskip}\usepackage[hidelinks]{hyperref}
\definecolor{tcetblue}{RGB}{0,0,0}\definecolor{tcetaccent}{RGB}{0,0,0}
\definecolor{codebg}{RGB}{245,246,248}\definecolor{codeframe}{RGB}{215,219,225}
\definecolor{codekw}{RGB}{0,0,0}\definecolor{codecomment}{RGB}{0,0,0}
\definecolor{codestr}{RGB}{0,0,0}\definecolor{boxframe}{RGB}{0,0,0}
\hypersetup{colorlinks=true,linkcolor=tcetblue,urlcolor=tcetaccent,citecolor=tcetblue,
  pdftitle={TCET Coding Platform -- Module 7: Deployment \& Future Scope},pdfauthor={Centre of Excellence, TCET}}
\lstdefinestyle{tcp}{backgroundcolor=\color{codebg},basicstyle=\ttfamily\footnotesize,
  keywordstyle=\color{codekw}\bfseries,commentstyle=\color{codecomment}\itshape,
  stringstyle=\color{codestr},numberstyle=\tiny\color{gray},frame=single,
  rulecolor=\color{codeframe},breaklines=true,showstringspaces=false,tabsize=2,
  columns=fullflexible,keepspaces=true,xleftmargin=0.6em,framexleftmargin=0.6em}
\lstset{style=tcp}
\newcommand{\code}[1]{\texttt{#1}}
\newenvironment{callout}[1]{\par\vspace{4pt}\noindent
  \begin{tabular}{@{}p{0.014\linewidth}@{\hspace{0.6em}}p{0.93\linewidth}@{}}
  \textcolor{boxframe}{\rule{2pt}{\baselineskip}} & \textbf{\textcolor{tcetblue}{#1}}\\[2pt] & }{\end{tabular}\par\vspace{6pt}}
\titleformat{\chapter}[display]{\bfseries\fontsize{18}{22}\selectfont}{\bfseries\fontsize{14}{17}\selectfont Chapter \thechapter}{0.2em}{}
\titlespacing*{\chapter}{0pt}{6pt}{22pt}
\titleformat{\section}{\bfseries\fontsize{16}{19}\selectfont}{\thesection}{0.6em}{}
\titleformat{\subsection}{\bfseries\fontsize{14}{17}\selectfont}{\thesubsection}{0.5em}{}
\pagestyle{fancy}\fancyhf{}\renewcommand{\headrulewidth}{0.4pt}
\fancyhead[L]{\small\itshape TCET Coding Platform}\fancyhead[R]{\small\itshape Module 7 \textbullet\ Deployment \& Future}
\fancyfoot[C]{\thepage}
\fancypagestyle{plain}{\fancyhf{}\fancyfoot[C]{\thepage}\renewcommand{\headrulewidth}{0pt}}
\newcommand{\note}[1]{\par\vspace{3pt}\noindent{\color{tcetaccent}\textbf{Note.}}~#1\par\vspace{3pt}}
\begin{document}
\begin{titlepage}\centering
\vspace*{1.4cm}{\color{tcetaccent}\rule{\linewidth}{1.2pt}}\\[0.5cm]
{\LARGE\bfseries TCET Coding Platform\par}
\vspace{0.4cm}{\Large Official Technical Documentation\par}{\large Draft~1 \textemdash\ Testing Phase\par}
{\color{tcetaccent}\rule{\linewidth}{1.2pt}}\\[1.4cm]
{\Huge\bfseries\color{tcetblue} Module 7\par}\vspace{0.3cm}
{\Huge Deployment, Infrastructure,\\[0.2cm] Operations \& Future Scope\par}\vspace{2.0cm}
{\large\bfseries Powered by\par}{\Large\color{tcetblue} Centre of Excellence (CoE), TCET\par}\vfill
\begin{tabular}{rl}\textbf{Module} & 7 of 7\\[3pt]
\textbf{Scope} & Topology, config, scaling, runbooks, roadmap\\[3pt]
\textbf{Status} & Draft~1 (living document until Stage~2)\\[3pt]
\textbf{Audience} & DevOps, Officials, Engineering\\\end{tabular}
\vspace{0.8cm}{\small\itshape \today\par}\end{titlepage}
\pagenumbering{roman}\tableofcontents\clearpage\pagenumbering{arabic}

% =====================================================================
\chapter{Deployment Topology}
% =====================================================================
\section{Scope of this Module}
This module specifies how the platform is deployed and operated: its network
topology, the complete configuration surface, scaling guidance, day-to-day
runbooks, the testing strategy, and the forward-looking roadmap that will guide
work beyond this draft.

\section{A Private-by-Default Deployment}
The production deployment is intentionally private: the application runs on a host
that is not publicly exposed, all public traffic is routed through a secure tunnel,
and administrative access is restricted to a private network.
\begin{longtable}{@{}p{0.26\linewidth} p{0.66\linewidth}@{}}
\toprule
\textbf{Element} & \textbf{Role}\\
\midrule
\endhead
Private admin network & Restricts administrative / SSH access; the host is not directly reachable from the internet.\\
Secure tunnel & Exposes local services as public hostnames via a reverse proxy.\\
Frontend hostname & Public entry point routed to the frontend.\\
API hostname & Public entry point routed to the backend API.\\
Process supervisor & Supervises the application processes and restarts them on failure.\\
Container runtime & Runs the execution (Judge0) stack.\\
\bottomrule
\end{longtable}

\section{Topology Diagram}
\begin{lstlisting}[language={}]
        Internet
           |
     Secure tunnel  (public hostnames -> local ports)
           |
     +-----+------------------------------+
     |  Host (private; admin via VPN)     |
     |                                    |
     |  [Frontend]   [API server]---------|--> MongoDB
     |                 |   \              |
     |                 |    \--> Redis <--|--- [Worker]
     |                 |                  |        |
     |                 +--> schedulers    |        +--> Judge0 (Docker)
     +------------------------------------+
\end{lstlisting}

\section{Trust Boundary Reminder}
The backend reads the proxy's forwarded headers and enforces the trusted-proxy
source-IP allowlist described in Module~3. The reverse proxy is responsible for
stripping any client-supplied identity headers so identity cannot be spoofed.

% =====================================================================
\chapter{Process Supervision \& Execution Stack}
% =====================================================================
\section{Supervised Processes}
\begin{longtable}{@{}p{0.28\linewidth} p{0.64\linewidth}@{}}
\toprule
\textbf{Process} & \textbf{Description}\\
\midrule
\endhead
API server & The Express API; the embedded worker is disabled here in production.\\
Submission worker & The dedicated queue consumer, isolating judging load from request handling.\\
Frontend & Serves the built single-page application.\\
\bottomrule
\end{longtable}
Each process runs with automatic restart and a memory ceiling, so a runaway
process is recycled without manual intervention.

\section{The Execution Stack}
The execution engine is deployed as a container stack with a sandbox override.
Helper scripts manage its lifecycle:
\begin{longtable}{@{}p{0.34\linewidth} p{0.58\linewidth}@{}}
\toprule
\textbf{Operation} & \textbf{Action}\\
\midrule
\endhead
Bring up & Start the local execution stack.\\
Bring down & Stop the stack.\\
Status & Report health / readiness.\\
Test sandbox & Verify sandboxed execution end to end.\\
Test languages & Verify the language runtime matrix.\\
\bottomrule
\end{longtable}

% =====================================================================
\chapter{Configuration Reference}
% =====================================================================
All backend configuration is validated at startup; the process refuses to start
on invalid configuration. The tables group the configuration surface by concern.

\section{Identity \& Security}
\begin{xltabular}{\linewidth}{@{}l X@{}}
\toprule
\textbf{Setting} & \textbf{Purpose}\\
\midrule
\endhead
SSO base URL & Base URL of the CoE SSO, used for cookie-domain logout.\\
Frontend base URL & Canonical frontend origin.\\
Token signing secret & Secret for verifying the CoE token; must meet a minimum strength.\\
Require trusted proxy & Enforce trusted-proxy source checks (kept on in production).\\
Trusted proxy allowlist & Comma-separated proxy source IPs / CIDRs.\\
CORS origins & Comma-separated allowed origins.\\
\bottomrule
\end{xltabular}

\section{Data, Queue \& Execution}
\begin{xltabular}{\linewidth}{@{}l X@{}}
\toprule
\textbf{Setting} & \textbf{Purpose}\\
\midrule
\endhead
Database URI \& name & MongoDB connection and database name.\\
Redis connection & Host, port, database, and password for the queue.\\
Queue name & The submission queue name.\\
Worker concurrency & Number of jobs a worker processes in parallel.\\
Embedded worker toggle & Whether the API hosts a worker (disabled in production).\\
Execution provider & Judge0 or the stub provider.\\
Judge0 endpoint & Base URL and optional cloud credentials.\\
Poll cadence & Result polling interval and timeout.\\
Sandbox pool size & Number of recycled isolate box identifiers leased per worker container (default 100; only needs to exceed worker concurrency). See Module~4.\\
Sandbox lock directory & Filesystem location of the per-box lock files used to lease sandbox identifiers.\\
\bottomrule
\end{xltabular}

\section{Scheduling \& Scoring}
\begin{xltabular}{\linewidth}{@{}l X@{}}
\toprule
\textbf{Setting} & \textbf{Purpose}\\
\midrule
\endhead
Stale threshold & Age after which a stuck submission is re-enqueued.\\
Finalizer interval & Cadence of the expired-attempt finalizer (zero disables it).\\
Default problem limits & Default CPU time and memory limits for problems.\\
Rating weights & Difficulty-weighted rating awards.\\
Environment \& port & Runtime mode and API port.\\
\bottomrule
\end{xltabular}

% =====================================================================
\chapter{Scaling Guidance}
% =====================================================================
\section{Principles}
\begin{itemize}[leftmargin=1.5em]
  \item \textbf{Absorb spikes with the queue.} Submissions are buffered, so a
        burst of many concurrent submits does not overwhelm execution; work
        drains at worker capacity.
  \item \textbf{Scale judging horizontally.} Run multiple worker processes or
        hosts against the same queue and tune concurrency to match the execution
        stack's throughput.
  \item \textbf{Separate API and worker.} Keep the worker out of the API process
        so request latency is unaffected by judging load.
  \item \textbf{Right-size limits.} Global and per-user rate limits and the
        payload cap protect the API under mass load.
\end{itemize}

\section{Capacity Planning}
\begin{longtable}{@{}p{0.32\linewidth} p{0.60\linewidth}@{}}
\toprule
\textbf{Dimension} & \textbf{Lever}\\
\midrule
\endhead
Submission throughput & Worker count $\times$ concurrency, bounded by execution capacity.\\
API request throughput & Horizontal API instances behind the proxy; rate limits.\\
Storage growth & Retention and archival policy (see Future Scope).\\
Execution capacity & Judge0 host sizing and language runtime availability.\\
\bottomrule
\end{longtable}

% =====================================================================
\chapter{Operational Runbooks}
% =====================================================================
\section{Local Development}
\begin{enumerate}[leftmargin=1.6em]
  \item Install and run the backend (and, optionally, a separate worker).
  \item Install and run the frontend, pointed at the backend base URL.
  \item Bring up the local execution stack and verify its status.
  \item Use the local mock SSO to obtain a development identity (never in
        production).
  \item Seed baseline data if needed.
\end{enumerate}

\section{Production Build \& Start}
\begin{enumerate}[leftmargin=1.6em]
  \item Build the backend and frontend.
  \item Start the API, worker, and frontend under the process supervisor.
  \item Ensure the datastore, queue, and execution stack are reachable.
\end{enumerate}

\section{Troubleshooting}
\begin{longtable}{@{}p{0.34\linewidth} p{0.58\linewidth}@{}}
\toprule
\textbf{Symptom} & \textbf{Likely cause / action}\\
\midrule
\endhead
Unauthorized source & The request source IP is not allowlisted; update the allowlist and restart.\\
Missing authentication & The request did not pass through the trusted auth path; verify the proxy injects identity.\\
Account not active & The upstream identity marks the user inactive; confirm the identity payload.\\
Failed to fetch & Check the API base URL, CORS origins, and ingress reachability.\\
Execution failures & Check execution-stack status and the sandbox/language self-tests; confirm the endpoint.\\
Submissions stuck & Stale recovery re-enqueues them; verify the worker and queue are healthy.\\
\bottomrule
\end{longtable}

\section{Health \& Observability}
A basic status route is public; health and database-diagnostic routes answer only
trusted internal sources and must not be exposed publicly. Security-relevant
authentication failures and background-scheduler outcomes are logged to the
process console for aggregation.

% =====================================================================
\chapter{Testing \& Quality Strategy}
% =====================================================================
\section{Test Levels}
\begin{longtable}{@{}p{0.30\linewidth} p{0.62\linewidth}@{}}
\toprule
\textbf{Level} & \textbf{Coverage}\\
\midrule
\endhead
Backend integration & The API run against in-memory repositories, a stub execution provider, and a fixed clock \textemdash\ covering auth, submissions, contest scoring, and recovery deterministically.\\
Execution \& wrapper & Validating the code wrapper and the execution provider's language resolution and verdict normalization.\\
Frontend & Client mapping, submission polling, and editor helpers.\\
Type safety & Whole-repository static type checking.\\
Load testing & Scripts that exercise the submission queue.\\
\bottomrule
\end{longtable}

\section{Why Testability was Designed In}
Because services depend on interfaces and an injected clock, the entire system
can be exercised deterministically without production infrastructure. This is a
first-class quality mechanism, not an afterthought: it makes regressions in
scoring, timing, and recovery catchable in automated tests.

% =====================================================================
\chapter{Future Scope \& Roadmap}
% =====================================================================
The following items are candidate work beyond Draft~1, recorded so the Stage~2
review can prioritize them. They are intentionally scoped, not committed. The
headline direction is to evolve the platform from an assessment platform into a
\textbf{full self-learning environment} \textemdash\ pairing guided, AI-assisted
practice with hands-on labs for programming, databases, and operating systems.

\section{AI-Assisted, Answer-Preserving Learning}
A central pillar of the roadmap is an AI layer that helps students \emph{learn}
rather than copy. Unlike platforms that reveal a canonical solution, this layer
never hands over the answer; it reasons about the student's own code.

\subsection{Solution Optimality Analysis}
After a student submits a working solution, the AI analyzes \emph{their} code and
reports how optimal it is \textemdash\ estimated time and space complexity,
redundant work, unnecessary allocations, and whether a better algorithmic class
exists (for example, ``your approach is quadratic; a log-linear solution is
possible''). It grades the solution's quality and explains \emph{why}, without
printing an optimal solution the student can simply paste back.

\subsection{Real-Time, Non-Revealing Hints}
While a student is coding, the AI analyzes the in-progress code and, when the
student is stuck or heading down a wrong path, offers \emph{graduated hints} that
nudge them toward the fix \textemdash\ pointing at the failing idea, the
off-by-one, or the missing case, escalating in specificity only as needed, but
always stopping short of writing the answer. This mirrors a good teaching
assistant: it unblocks without solving.

\subsection{Guardrails}
The AI operates under strict integrity rules: it is disabled (or restricted to
complexity feedback only) during live, rated contests, and its hint verbosity is
configurable by faculty per problem or per course, so it can be tuned from a
gentle nudge to off.

\begin{callout}{Design intent: answer-preserving}
The AI's objective is the student's understanding, not task completion. It
critiques and hints on the student's own work and never substitutes a model
solution for it \textemdash\ so it accelerates learning without becoming a
shortcut to the answer.
\end{callout}

\section{Integrated Learning \& Practice Environments}
Beyond competitive-style problems, the platform will host structured, hands-on
learning tracks so a student can practice an entire course inside the platform.

\subsection{TCET Practical Labs}
Digitize the institution's practical laboratory sessions: each lab becomes a
structured, auto-evaluated assignment with objectives, starter scaffolding,
sandboxed execution, and automatic capture of student work \textemdash\ replacing
manual lab submissions and giving faculty a single dashboard for practical-exam
conduct and assessment. This directly targets the routine laboratory workflow
that today is spread across machines and paper.

\subsection{DBMS \& SQL Learning Module}
An interactive database track where students write and run real SQL against a
sandboxed database engine, with schema-and-query exercises, result-set checking,
and the same AI optimality feedback applied to queries (index usage, join
strategy, query cost). This extends the platform from programming into core
database coursework.

\subsection{Linux / Operating-Systems Terminal Practice}
A browser-based, sandboxed Linux terminal environment where students practice
shell commands, file-system navigation, permissions, processes, and scripting
\textemdash\ a safe operating-systems laboratory, with task-based challenges
(``create this directory tree,'' ``find and stop this process'') that are
automatically verified against the container's state.

\begin{callout}{One spine, many environments}
These environments reuse the platform's existing spine \textemdash\ sandboxed
execution, per-attempt state, proctoring, and analytics \textemdash\ so labs, SQL
exercises, and terminal challenges are authored, assigned, proctored, and graded
through the same faculty tooling documented in this record.
\end{callout}

\section{Assessment \& Integrity Enhancements}
\begin{itemize}[leftmargin=1.5em]
  \item Optional integrations for invigilated / lockdown environments and, where
        policy permits, webcam-assisted proctoring, to close the browser-only
        limitations noted in Module~6.
  \item Plagiarism and code-similarity detection across contest submissions.
  \item Partial-credit and subtask-based grading for coding questions, and
        configurable penalty policies.
  \item Randomized question sets per student to reduce answer sharing.
\end{itemize}

\section{Platform Capabilities}
\begin{itemize}[leftmargin=1.5em]
  \item An explicit administrator tier distinct from faculty, with
        institution-wide oversight and audit dashboards.
  \item Richer analytics: cohort comparisons, topic-wise mastery, and
        outcome-based reporting for accreditation.
  \item Real-time submission status via push channels to replace client polling.
  \item Additional language runtimes and per-department language policies.
\end{itemize}

\section{Operations \& Reliability}
\begin{itemize}[leftmargin=1.5em]
  \item Centralized structured logging, metrics, and alerting.
  \item Explicit database indexing and archival strategy for large historical
        datasets, with a documented backup / restore procedure.
  \item Horizontal auto-scaling of judging workers driven by queue depth.
  \item A migration to retire the legacy repository naming for clarity.
\end{itemize}

\section{Documentation \& Governance}
\begin{itemize}[leftmargin=1.5em]
  \item Promotion of this Draft~1 to a Stage~2 baseline with formal change
        control and an API reference appendix.
  \item Data-retention, privacy, and consent documentation aligned to
        institutional policy.
\end{itemize}

\section{Indicative Roadmap Phasing}
\begin{longtable}{@{}p{0.20\linewidth} p{0.72\linewidth}@{}}
\toprule
\textbf{Horizon} & \textbf{Candidate focus}\\
\midrule
\endhead
Near term & Practical Labs digitization; structured logging and metrics; indexing and backups.\\
Mid term & AI optimality analysis and real-time hints; DBMS/SQL module; plagiarism detection.\\
Longer term & Linux terminal environment; administrator tier and audit dashboards; auto-scaling; accreditation analytics.\\
\bottomrule
\end{longtable}

% =====================================================================
\chapter{Conclusion}
% =====================================================================
\section{Where the Platform Stands}
As documented across these seven modules, the TCET Coding Platform is a
functioning, institution-owned coding and assessment platform engineered for
scale, integrity, and ownership: federated identity behind a trusted proxy, an
asynchronous sandboxed execution pipeline, a rigorous contest engine with
deferred grading and a feedback gate, a layered proctoring subsystem, and a
private, supervised deployment.

\section{The Path Forward}
The roadmap turns this foundation toward a comprehensive self-learning
environment \textemdash\ AI that helps students reason about their own code
without revealing answers, and hands-on labs for programming, databases, and
operating systems \textemdash\ all built on the same spine of sandboxed
execution, integrity, and analytics. This document, Draft~1, is the baseline
against which that progress will be reviewed at Stage~2.

\vspace{0.8cm}
\begin{center}
\textit{--- End of Module 7, and of the Draft~1 documentation set. Subject to revision at Stage~2. ---}
\end{center}
\end{document}
